\section{Equations with a model reference}
\label{sec:general-measure-relative}

We now prove \cref{thm:als-general-measure}. Write the given measure as $\mu=f\nu$, where $\nu$ is the Monge--Amp\`ere measure of a bounded potential $\psi$. To keep the reference $\varphi$, we first smooth $\psi$: the prescribed ordinary-type theorem applies to the approximating data, exponential comparison provides a common lower envelope, and moving-measure continuity preserves the equation in the limit. With $\psi$ fixed, we then approximate the integrable density $f$ by bounded densities. Here the solutions for the reference zero give upper bounds, while positive lower mass prevents collapse.

\begin{proof}[of \cref{thm:als-general-measure}]
The proof has four steps: continuous positive densities, bounded densities, integrable densities, and finally the zero exponent.

Let $\mu$ be the given nonzero finite non-pluripolar measure. Apply \cref{lem:bounded-measure-carrier} to the Hermitian, hence semipositive and big, background $\omega$ to write
\[
 \mu=f\nu,\qquad \nu=\omega_\psi^n,\qquad
 \psi\in\PSH(X,\omega)\cap L^\infty(X),\qquad
 0\leq f\in L^1(\nu).
\]
\emph{Step 1: positive continuous density.} Fix $\lambda>0$ and suppose first that $f$ is continuous and strictly positive. By \cref{thm:bounded-exact-regularization}, choose smooth strictly $\omega$-plurisubharmonic $\psi_j\downarrow\psi$. These potentials are uniformly bounded; write $\|\psi_j\|_\infty\leq L$ and $\nu_j=\omega_{\psi_j}^n$. Bounded Bedford--Taylor continuity gives $\nu_j\rightharpoonup\nu$. The measure $f\nu_j$ has a continuous positive density with respect to smooth volume, so it belongs to every finite $L^p$ space. Hence the prescribed ordinary-type theorem \cref{thm:als-prescribed-lp} gives $u_j\sim\varphi$ with
\[
 \mu_{u_j}=\mathrm{e}^{\lambda u_j}f\nu_j.
\]

We first bound the $u_j$ from above, uniformly in $j$. Choose $K$ with
$\mathrm{e}^{-\lambda(\psi_j+K)}\leq\min_Xf$ for all $j$. Then
\[
 \mathrm{e}^{-\lambda(\psi_j+K)}\omega_{\psi_j+K}^n
 \leq f\nu_j
 =\mathrm{e}^{-\lambda u_j}\mu_{u_j}
 \leq \mathrm{e}^{-\lambda u_j}\mathrm{MA}^{\mathrm{ext}}_\omega(u_j),
\]
where the last inequality follows from $\omega\geq\theta$ and positivity of extended mixed products. Apply \cref{cor:singular-exponential-order} in the background $\omega$, with first potential $\psi_j+K$ and second potential $u_j$. Its hypotheses hold: the first potential is bounded, so it admits a smaller-background envelope, and $u_j\preceq\psi_j+K$. Thus $u_j\leq\psi_j+K\leq L+K=:C$.

By \eqref{eq:relative-normalized-model-bound}, $u_j\leq\varphi+C$. Since
\[
 \mu_{u_j}\leq\|f\|_\infty \mathrm{e}^{\lambda u_j}\omega_{\psi_j}^n,
\]
\cref{prop:exponential-common-rooftop} gives a single $h\in\mathcal E_\varphi$ below all $u_j$. The bounds $h\leq u_j\leq\varphi+C$ and \cref{thm:normalized-compactness} give, after extraction, an $L^1$ limit $u\not\equiv-\infty$. The moving-measure theorem \cref{thm:relative-moving-measure-continuity} now applies, because $\psi_j\to\psi$ in capacity and $\mu_{u_j}\leq \mathrm{e}^{\lambda C}\|f\|_\infty\nu_j$. It gives $u\in\mathcal E_\varphi$, convergence in capacity and $\mu_{u_j}\rightharpoonup\mu_u$.

The right-hand sides converge weakly to $\mathrm{e}^{\lambda u}f\nu$. Indeed, $\mathrm{e}^{\lambda u_j}-\mathrm{e}^{\lambda u}\to0$ in capacity with a common bound, so its integral against $\nu_j$ tends to zero by the uniform capacity estimate for these measures. The fixed function $\mathrm{e}^{\lambda u}$ is bounded and quasi-continuous. Approximating it by a continuous function outside an open set of arbitrarily small capacity, bounded product continuity and the same capacity estimate allow this fixed weight in $\nu_j\rightharpoonup\nu$. Multiplication by the continuous function $f$ and by a continuous test function preserves both estimates. Hence $\mu_u=\mathrm{e}^{\lambda u}f\nu$.

\emph{Step 2: bounded density.} Suppose next that $0\leq f\in L^\infty(\nu)$ and $\int_Xf\,\mathrm{d}\nu>0$. Choose positive continuous functions $f_j$, uniformly bounded by a constant $F$, with $f_j\to f$ almost everywhere for $\nu$. Such a sequence is obtained from the density of continuous functions in $L^1(\nu)$ by clipping to a common interval, adding positive constants tending to zero, and extracting a summably convergent subsequence.

There is a bounded nonnegative function $g$ with $g\nu\ne0$ and $g\leq f_j$ for all $j$. Indeed, choose $a>0$ with $\nu(\{f\geq2a\})>0$. By Egorov's theorem, there is a measurable subset $E\subseteq\{f\geq2a\}$ of positive measure on which the convergence is uniform; hence $f_j\geq a$ on $E$ for all $j\geq N$, for some $N$. The finitely many remaining functions $f_1,\ldots,f_{N-1}$ are positive and continuous, so they have positive minima on $X$. If $b>0$ is a lower bound for these minima and for $a$, then $g=b\mathbf1_E$ works.

By Step 1, let $u_j\in\mathcal E_\varphi$ solve $\mu_{u_j}=\mathrm{e}^{\lambda u_j}f_j\nu$. The theorem for the reference zero, \cref{thm:general-measure-minimal}, applied to the finite nonzero measure $g\nu$, gives $s\in\mathcal E_0$ satisfying
\[
 \mathrm{MA}^{\mathrm{ext}}_\theta(s)=\mathrm{e}^{\lambda s}g\nu.
\]
Apply \cref{cor:relative-extended-exponential-order} with reference zero, first potential $s\in\mathcal E_0$ and second potential $u_j\preceq0$. Since $g\nu\leq f_j\nu$, it gives $u_j\leq s$. With $C=\sup_Xs$ we have
\[
 \mu_{u_j}\leq \mathrm{e}^{\lambda C}F\nu,
 \qquad
 \kappa\leq \mathrm{e}^{\lambda\sup_Xu_j}F\nu(X),
\]
where the second inequality uses $\mu_{u_j}(X)\geq\kappa$. These two bounds on the suprema give a subsequence converging in $L^1$ to some $u\not\equiv-\infty$. The fixed-measure continuity theorem \cref{thm:relative-fixed-measure-continuity} puts $u$ in $\mathcal E_\varphi$ and identifies the weak limit of $\mu_{u_j}$ as $\mu_u$.

By \cref{lem:fixed-npp-integrals}, $\mathrm{e}^{\lambda u_j}\to \mathrm{e}^{\lambda u}$ in $L^1(\nu)$, and $f_j\to f$ in $L^1(\nu)$ by dominated convergence. Since both factors are uniformly bounded,
\[
 \mathrm{e}^{\lambda u_j}f_j\longrightarrow \mathrm{e}^{\lambda u}f
 \quad\hbox{in }L^1(\nu).
\]
Thus the right-hand sides converge in total variation, which proves the equation for $f$.

\emph{Step 3: integrable density.} For general $f\in L^1(\nu)$, apply Step 2 to $f_j=\min(f,j)$ and obtain $u_j\in\mathcal E_\varphi$ with $\mu_{u_j}=\mathrm{e}^{\lambda u_j}f_j\nu$. Since $f_j\leq f_{j+1}$, exponential order gives $u_{j+1}\leq u_j$. Furthermore
\[
 \kappa\leq \mathrm{e}^{\lambda\sup_Xu_j}\mu(X),
 \qquad \mu_{u_j}\leq \mathrm{e}^{\lambda\sup_Xu_1}\mu.
\]
Hence the decreasing limit $u$ is not identically $-\infty$, and the fixed-measure continuity theorem gives $u\in\mathcal E_\varphi$ and $\mu_{u_j}\rightharpoonup\mu_u$. Dominated convergence gives $\mathrm{e}^{\lambda u_j}f_j\to \mathrm{e}^{\lambda u}f$ in $L^1(\nu)$, so $\mu_u=\mathrm{e}^{\lambda u}\mu$. Uniqueness follows from \cref{cor:relative-extended-exponential-order}, applied in both directions.

\emph{Step 4: zero exponent.} For each positive integer $j$, Step 3 with exponent $1/j$ gives $u_j\in\mathcal E_\varphi$ with $\mu_{u_j}=\mathrm{e}^{u_j/j}\mu$. The zero-exponent part of the proof of \cref{thm:general-measure-minimal} applies to these $u_j$, with the same comparison potential $s\in\mathcal E_0$, the solution of $\mathrm{MA}^{\mathrm{ext}}_\theta(s)=\mathrm{e}^s\mu$. Indeed, the comparison $u_j\leq s+(j-1)\max(0,\sup_Xs)$ made there uses exponential order with reference zero, first potential in $\mathcal E_0$ and second potential $u_j\preceq0$; the lower mass bound $v_-(\theta)$ is replaced by $\kappa$; and \cref{thm:relative-fixed-measure-continuity} is applied with reference $\varphi$. We obtain $v\in\mathcal E_\varphi$ with $\sup_Xv=0$ and $c>0$ with $\mu_v=c\mu$. Finally, the scalar assertion of \cref{cor:relative-extended-exponential-order}, applied in both directions, shows that $c$ is unique.
\end{proof}
